% AA-72
\chapter{Funkcja partycji i wzór krotności Kostanta}
\label{aa:ch:kostant-multiplicity}

W AA-70 uzyskaliśmy wzór na cały charakter reprezentacji, a w AA-71 na jej wymiar. Chcemy teraz obliczać krotność pojedynczej wagi bez uprzedniego wyznaczania wszystkich przestrzeni wagowych. W tym celu rozłożymy mianownik Weyla na czynniki i zinterpretujemy ich odwrotności kombinatorycznie.

Zachowujemy założenia AA-70. Niech $K$ będzie algebraicznie domkniętym ciałem charakterystyki zero, $\mathfrak g$ niezerową skończenie wymiarową półprostą algebrą Liego nad $K$, $\mathfrak h$ maksymalną podalgebrą toralną, $\Phi^+$ zbiorem dodatnich korzeni względem bazy $\Delta=(\alpha_1,\ldots,\alpha_r)$, $P$ kratą wag, $P_+$ zbiorem wag dominujących, $W$ grupą Weyla oraz $\rho=\frac12\sum_{\alpha\in\Phi^+}\alpha$. Oznaczamy przez $Q=\sum_i\mathbb Z\alpha_i$ kratę korzeni, a przez $Q_+=\sum_i\mathbb Z_{\ge0}\alpha_i$ jej nieujemny stożek. Dla $\gamma=\sum_i n_i\alpha_i\in Q_+$ definiujemy wysokość $\operatorname{ht}\gamma=\sum_i n_i$.

\section{Funkcja partycji korzeni}

\begin{definicja}
\label{aa:def:kostant-partition}
Funkcją partycji Kostanta nazywamy odwzorowanie $\mathcal P:P\to\mathbb Z_{\ge0}$ dane wzorem
\[
 \mathcal P(\gamma)=
 \#\left\{(k_\alpha)_{\alpha\in\Phi^+}\in
 \mathbb Z_{\ge0}^{\Phi^+}:
 \sum_{\alpha\in\Phi^+}k_\alpha\alpha=\gamma\right\}.
\]
Liczymy rodziny krotności korzeni, nie uporządkowane ciągi składników. W szczególności $\mathcal P(\gamma)=0$ dla $\gamma\notin Q_+$, a $\mathcal P(0)=1$.
\end{definicja}
Liczba ta jest skończona. Istotnie, każdy dodatni korzeń należy do $Q_+$ i ma wysokość co najmniej jeden. Jeżeli $\gamma\in Q_+$, to z równości w definicji wynika
$\sum_\alpha k_\alpha\operatorname{ht}\alpha=\operatorname{ht}\gamma$,
więc każde $k_\alpha$ jest nie większe niż $\operatorname{ht}\gamma$. Dla $\gamma=0$ wszystkie $k_\alpha$ muszą być zerowe, co uzasadnia podaną wartość.

Ustalmy dowolne uporządkowanie $\Phi^+=(\beta_1,\ldots,\beta_N)$, gdzie $N=|\Phi^+|$. Przez $\mathcal P_j(\gamma)$ oznaczamy analogiczną liczbę rozkładów z wykorzystaniem wyłącznie $\beta_1,\ldots,\beta_j$, dla $0\le j\le N$. Zatem $\mathcal P_0(\gamma)=1$ dla $\gamma=0$ i zero w przeciwnym razie, a $\mathcal P_N=\mathcal P$.

\begin{lemat}[Rekurencja funkcji partycji]
\label{aa:lem:kostant-partition-recursion}
Dla $1\le j\le N$ i $\gamma\in P$ zachodzi
\[
 \mathcal P_j(\gamma)=
 \mathcal P_{j-1}(\gamma)+\mathcal P_j(\gamma-\beta_j).
\]
\end{lemat}
\begin{proof}
Dzielimy rozkłady liczone przez $\mathcal P_j(\gamma)$ na dwie rozłączne klasy. W pierwszej krotność $\beta_j$ jest zerowa, więc zapomnienie tej współrzędnej daje bijekcję z rozkładami liczonymi przez $\mathcal P_{j-1}(\gamma)$. W drugiej krotność $\beta_j$ jest dodatnia. Zmniejszenie jej o jeden daje rozkład $\gamma-\beta_j$ przy użyciu pierwszych $j$ korzeni. Odwrotnie, zwiększenie tej współrzędnej o jeden odzyskuje dokładnie jeden rozkład z drugiej klasy. Obie bijekcje działają także wtedy, gdy odpowiedni zbiór jest pusty. Dodanie liczności daje tezę.
\end{proof}

\section{Funkcja tworząca i formalne mnożenie}

Nie wolno odwracać $1-e^{-\alpha}$ w samym pierścieniu $\mathbb Z[P]$: geometryczny szereg ma na ogół nieskończony nośnik. Wystarczy jednak określić działanie skończonych elementów pierścienia na rodzinach współczynników.

Niech $\mathscr F$ będzie grupą abelową wszystkich funkcji $f:P\to\mathbb Z$ z dodawaniem punktowym. Zapisujemy formalnie $f=\sum_{\nu\in P}f(\nu)e^\nu$, bez twierdzenia o zbieżności tej sumy. Dla $H=\sum_{\eta\in P}h_\eta e^\eta\in\mathbb Z[P]$ określamy
\[
 (Hf)(\nu)=\sum_{\eta\in P}h_\eta f(\nu-\eta).
\]
Suma jest skończona, gdyż $H$ ma skończony nośnik. Dla $H,G\in\mathbb Z[P]$ i $f\in\mathscr F$ obie strony równości $(HG)f=H(Gf)$ przy $\nu$ są sumą
$\sum_{\eta,\xi}h_\eta g_\xi f(\nu-\eta-\xi)$;
sumujemy po skończonych nośnikach $H,G$. Ponadto $e^0f=f$, a rozdzielność wynika z rozdzielności dodawania i mnożenia liczb całkowitych. Otrzymujemy więc moduł nad $\mathbb Z[P]$, zawierający ten pierścień przez identyfikację elementów o skończonym nośniku z ich funkcjami współczynników. Nie definiujemy mnożenia dowolnych dwóch funkcji z $\mathscr F$.

\begin{lemat}[Funkcja tworząca]
\label{aa:lem:kostant-generating}
Dla $0\le j\le N$ określmy $B_j\in\mathscr F$ przez $B_j(\nu)=\mathcal P_j(-\nu)$. Wtedy
\[
 \left(\prod_{i=1}^j(1-e^{-\beta_i})\right)B_j=e^0.
\]
W szczególności symboliczny zapis funkcji tworzącej
\[
 B_N=\sum_{\gamma\in Q_+}\mathcal P(\gamma)e^{-\gamma}
     =\prod_{\alpha\in\Phi^+}\left(\sum_{k\ge0}e^{-k\alpha}\right)
\]
jest poprawny współczynnik po współczynniku.
\end{lemat}
\begin{proof}
Z definicji działania i Lematu~\ref{aa:lem:kostant-partition-recursion} dla $j\ge1$ mamy
\[
 \begin{aligned}
 ((1-e^{-\beta_j})B_j)(\nu)
 &=\mathcal P_j(-\nu)-\mathcal P_j(-\nu-\beta_j)\\
 &=\mathcal P_{j-1}(-\nu)=B_{j-1}(\nu).
 \end{aligned}
\]
Ponieważ $B_0=e^0$, kolejne zastosowanie tych równości dla $j,j-1,\ldots,1$ wraz z łącznością działania daje pierwszą tezę.

W symbolicznym iloczynie szeregów wybór jednego składnika z każdego czynnika wyznacza rodzinę $(k_\alpha)$ i monom $e^{-\sum_\alpha k_\alpha\alpha}$. Dla ustalonego wykładnika $-\gamma$ istnieje tylko skończenie wiele takich rodzin, jak wykazano przy definicji. Współczynnik przy tym monomie jest dokładnie $\mathcal P(\gamma)$, a poza $-Q_+$ jest zerowy. Tym samym ten konkretny iloczyn nieskończonych szeregów jest określony współczynnik po współczynniku; nie zakładamy istnienia ogólnego iloczynu w $\mathscr F$.
\end{proof}

\section{Wzór krotności}

\begin{twierdzenie}[Wzór krotności Kostanta]
\label{aa:tw:kostant-multiplicity}
Przy powyższych założeniach dla każdej $\lambda\in P_+$ i $\mu\in P$ zachodzi
\[
 m_\lambda(\mu):=\dim_K L(\lambda)_\mu
 =\sum_{w\in W}\epsilon(w)\,
 \mathcal P\bigl(w(\lambda+\rho)-(\mu+\rho)\bigr),
 \qquad \epsilon(w)=\det_E w.
\]
Dla wag niewystępujących w reprezentacji lewa strona jest z definicji zerowa.
\end{twierdzenie}
\begin{proof}
Niech $F=\operatorname{ch}L(\lambda)\in\mathbb Z[P]$, $B=B_N$ oraz $C=\prod_{\alpha\in\Phi^+}(1-e^{-\alpha})\in\mathbb Z[P]$. Twierdzenie~\ref{aa:tw:weyl-character-formula} wraz z tożsamością mianownika AA-69 daje w pierścieniu $\mathbb Z[P]$
\[
 CF=e^{-\rho}A_{\lambda+\rho}
   =\sum_{w\in W}\epsilon(w)e^{w(\lambda+\rho)-\rho}.
\]
Z Lematu~\ref{aa:lem:kostant-generating} mamy $CB=e^0$ w $\mathscr F$. Działając na $B$ przez obie strony równości pierścieniowej, uzyskujemy
\[
 F=F(CB)=(CF)B
   =\sum_{w\in W}\epsilon(w)e^{w(\lambda+\rho)-\rho}B.
\]
W pierwszej równości identyfikujemy $F$ z funkcją współczynników; działanie $F$ na $e^0$ daje dokładnie tę funkcję. Druga równość wynika z łączności działania i przemienności pierścienia grupowego. Ostatnia suma jest skończona, bo $W$ jest skończona.

Współczynnik przy $e^\mu$ w składniku $e^{w(\lambda+\rho)-\rho}B$ jest równy
$B(\mu-w(\lambda+\rho)+\rho)=
\mathcal P(w(\lambda+\rho)-(\mu+\rho))$.
Współczynnik lewej strony to $m_\lambda(\mu)$ z definicji charakteru. Porównanie tych współczynników daje tezę dla każdej $\mu\in P$.
\end{proof}

\begin{wniosek}
\label{aa:cor:kostant-coset}
Jeżeli $\lambda\in P_+$ i $\mu\in P$ oraz $\lambda-\mu\notin Q$, to $m_\lambda(\mu)=0$.
\end{wniosek}
\begin{proof}
Dla każdego $\nu\in P$ odbicie proste spełnia
$s_i\nu-\nu=-\nu(h_i)\alpha_i\in Q$. Ponieważ odbicia proste generują $W$ i zachowują $P$, iterowanie daje $w\nu-\nu\in Q$ dla wszystkich $w\in W$. Dla $\nu=\lambda+\rho$ argument funkcji partycji we wzorze Kostanta różni się od $\lambda-\mu$ o element $Q$. Nie należy więc do $Q$, a tym bardziej do $Q_+$. Wszystkie składniki sumy znikają.
\end{proof}

\section{Algorytm obliczania krotności}

Dla ustalonej $\gamma\in P$ najpierw zapisujemy ją we współrzędnych bazy korzeni prostych. Jeżeli co najmniej jedna współrzędna nie jest nieujemną liczbą całkowitą, zwracamy zero. W przeciwnym razie niech $\gamma=\sum_i n_i\alpha_i$ i oznaczmy przez $b_j\in\mathbb Z_{\ge0}^r$ wektor współrzędnych $\beta_j$. Wystarczy tablica o indeksach $u=(u_1,\ldots,u_r)$ z $0\le u_i\le n_i$. Zaczynamy od tablicy $T_0(0)=1$, $T_0(u)=0$ dla $u\ne0$. Dla $j=1,\ldots,N$ wykonujemy kroki
\[
 T_j(u)=T_{j-1}(u)+
 \begin{cases}
 T_j(u-b_j),&u_i\ge(b_j)_i\text{ dla wszystkich }i,\\
 0,&\text{w przeciwnym razie}.
 \end{cases}
\]
Indeksy $u$ przetwarzamy według niemalejącej sumy współrzędnych. Ponieważ $b_j\ne0$ i ma nieujemne współrzędne, indeks $u-b_j$, jeśli należy do tablicy, ma mniejszą sumę współrzędnych i jego wartość została już obliczona. Indukcja po $j$ oraz po tej sumie, z zastosowaniem Lematu~\ref{aa:lem:kostant-partition-recursion}, daje
$T_j(u)=\mathcal P_j(\sum_i u_i\alpha_i)$.
Wynik to $T_N(n)$. Dwie tablice zajmują $O(\prod_i(n_i+1))$ komórek, a liczba dodawań całkowitoliczbowych nie przekracza $N\prod_i(n_i+1)$. Nie jest to oszacowanie złożoności bitowej: rozmiar liczb całkowitych może rosnąć. Kolejność indeksów można uzyskać również przez rosnący porządek leksykograficzny, ponieważ odjęcie niezerowego nieujemnego wektora zmniejsza indeks w tym porządku.

Do obliczenia $m_\lambda(\mu)$ wyznaczamy orbitę $\lambda+\rho$ przez kolejne odbicia proste, zapisując przy każdym elemencie znak iloczynu odbić. Waga $\lambda+\rho$ jest ściśle dominująca, więc regularna (jej parowania ze wszystkimi dodatnimi korzeniami są dodatnie). Z Lematu~\ref{aa:lem:weyl-regular-orbits} stabilizator jest trywialny, zatem elementy orbity odpowiadają bijektywnie elementom $W$, a przypisany znak nie zależy od wybranej drogi w orbicie. Skończoność $W$ zapewnia zakończenie enumeracji. Dla każdej wagi orbity wyznaczamy argument funkcji partycji, obliczamy jego wartość powyższą rekurencją i dodajemy ze znakiem. Wzór Kostanta gwarantuje poprawność wyniku; poszczególne sumy pośrednie mogą być ujemne, choć wynik końcowy jest nieujemny.

\section{Przykłady}

\begin{przyklad}[Typ $A_1$]
\label{aa:prz:kostant-a1}
Dla jedynego dodatniego korzenia $\alpha=2\omega$ mamy $\mathcal P(\gamma)=1$ dokładnie wtedy, gdy $\gamma$ jest nieujemną całkowitą wielokrotnością $2\omega$. Dla $\lambda=a\omega$, $a\ge0$, $\mu=b\omega$, $b\in\mathbb Z$, wzór daje
\[
 m_{a\omega}(b\omega)=\mathcal P((a-b)\omega)
                     -\mathcal P((-a-b-2)\omega).
\]
Jeżeli $a-b$ jest nieparzyste, oba argumenty mają nieparzyste współczynniki przy $\omega$ i wynik jest zerowy. Jeżeli $a-b$ jest parzyste, pierwszy składnik jest jeden dla $b\le a$, a drugi dla $b\le-a-2$. Ponieważ wtedy $b$ i $a$ mają tę samą parzystość, różnica jest jeden dokładnie dla $-a\le b\le a$ i zero w pozostałych przypadkach. Odtwarzamy krotności z AA-68.
\end{przyklad}

\begin{przyklad}[Typ $A_2$]
\label{aa:prz:kostant-a2}
Dodatnie korzenie to $\alpha_1,\alpha_2,\alpha_1+\alpha_2$. Dla $p,q\in\mathbb Z_{\ge0}$ rozkład $p\alpha_1+q\alpha_2$ opisują równania $k_1+k_3=p$ oraz $k_2+k_3=q$. Każdy wybór $0\le k_3\le\min(p,q)$ daje dokładnie jedną parę nieujemnych krotności $k_1=p-k_3$, $k_2=q-k_3$. Zatem
$\mathcal P(p\alpha_1+q\alpha_2)=1+\min(p,q)$.

Dla reprezentacji dołączonej $\lambda=\rho=\alpha_1+\alpha_2$ obliczmy $m_\rho(0)$. W bazie korzeni prostych orbita $2\rho$ składa się z wektorów
$(2,2),(0,2),(2,0),(-2,0),(0,-2),(-2,-2)$,
otrzymanych kolejno przez identyczność, dwa odbicia proste, dwa iloczyny dwóch odbić i iloczyn trzech odbić. Po odjęciu $\rho=(1,1)$ tylko pierwszy wektor pozostaje w $Q_+$; jego wkład wynosi $\mathcal P(\rho)=2$. Stąd $m_\rho(0)=2$.

Dla modułu trywialnego $L(0)$ i wagi $\mu=-\rho$ argumenty funkcji partycji to po prostu $w\rho$. W orbicie $\rho$ wektory o nieujemnych współrzędnych to $(1,1),(0,1),(1,0)$, o znakach odpowiednio $+1,-1,-1$. Pozostałe trzy wektory mają ujemne współrzędne i wnoszą zero. Otrzymujemy
$m_0(-\rho)=\mathcal P(\rho)-\mathcal P(\alpha_2)-\mathcal P(\alpha_1)=2-1-1=0$.
Przykład pokazuje, że sama liczba rozkładów $\lambda-\mu$ nie jest na ogół krotnością wagi: istotne jest znoszenie składników ze znakami.
\end{przyklad}

Wzór Kostanta wyznacza krotności poszczególnych wag za pomocą skończonych obliczeń. Nie jest to wzór na krotności składników nieprzywiedlnych w rozkładzie iloczynu tensorowego; takie krotności wymagają osobnego argumentu.
